%!TEX program = lualatex
\documentclass[ngerman,headheight=70pt]{scrartcl}
\usepackage[ngerman]{babel}
\usepackage{blindtext}
\usepackage[utf8]{luainputenc}
\usepackage{scrlayer-scrpage}
\usepackage[T1]{fontenc}
\usepackage{fontspec}
\usepackage{multicol}
\usepackage{enumitem}
\usepackage[german=quotes]{csquotes}
\usepackage[bottom=2cm,footskip=8mm]{geometry}
\setmainfont{Source Sans Pro}
\MakeOuterQuote{"}

\setlength{\parindent}{0pt}
\setlength{\parskip}{6pt}

\pagestyle{plain.scrheadings}
\thispagestyle{scrheadings}
\lohead{\vspace*{-1.5cm}\includegraphics[scale=0.4]{../images/stupa_logo_rot.pdf}}
\rohead{\\Präsidium des\\Studierendenparlaments}

\lofoot[Universität Hamburg · Präsidium des Studierendenparlaments\\\hrule Von-Melle-Park 5 · D-20146 Hamburg]{Universität Hamburg · Präsidium des Studierendenparlaments\\\hrule Von-Melle-Park 5 · D-20146 Hamburg}
\cofoot[]{}
\cohead[\thepage]{}

\renewcommand{\thesubsection}{TOP \arabic{subsection}}
\renewcommand{\thesubsubsection}{\arabic{subsubsection}.}
\setcounter{subsection}{-1}

\begin{document}
    UHH · StuPa-Präsidium · Von-Melle-Park 5 · D-20146 Hamburg

    %%% PAGE ONE (EDIT LATER) %%%%
    % Hier den Titel des Protokolls eintragen
    \section*{Protokoll der 8. Sitzung des Studierendenparlaments vom 27. Oktober 2016}

    % Protokollanten eintragen
    \textbf{Protokoll: Ramon Weilinger}\\
    \textbf{Ort: VMP5 0079}\\
    \textbf{Beginn: 18.29 Uhr}\\
    \textbf{Ende: 20.39 Uhr}

    % Wohl die komplizierteste Stelle im Protokoll
    % Als Faustregel sollten nicht mehr als 3 Namen auf einer Zeile stehen
    % Einfach an dieser Vorlage orientieren
    % Im Idealfall ist hier nicht viel zu ändern, da kaum jemand fehlt.
    % Rücktritte müssen natürlich pro Protokoll angepasst werden
    \vspace{0.5cm}
    \begin{tabular}{ll}
        Anwesend: & \\
            RCDS (5 Sitze): & Ramon Weilinger, Ramin Shakiba, Jennifer Maack, Benjamin Welling \\
             CampusGrün (14 Sitze): & Laura Franzen, Elena Rysikova, Philipp Droll, \\
                                   & Yasemin Günther, Mario Moldenhauer, Christian Lagod, \\
                                   & Jim Martens, Svenja Horn, Martin Sievert, \\
                                   & Armin Günther \\
             BGZM-Liste (2 Sitze) : Karen Martirosian & \\
             WiWi (2 Sitze): & \\
             Unicorns (5 Sitze): & Katharina Kucza, Johannes Peplow, Marielle Hermstrüwer \\
             Liste LINKS (3 Sitze): & Gunhild Berdal, Till Petersen, Sinah Mielich \\
             HWP (2 Sitze): &  \\
             MIN (4 Sitze): & \\
             SDS* (3 Sitze): & Mena Winkler, Jacob Petersein, Artur Brückmann \\
             Bart-LISTE (2 Sitze): & Timo Zeimet \\
             LHG (1 Sitz): & Tobias Heisig \\
             harte zeiten (1 Sitz): & Tobias Berking \\
             Jura (1 Sitz): & \\
             AL (2 Sitze): & \\
            & \\
        Entschuldigt: & Antonia Niecke (RCDS), Geoffrey Youett, Melf Johannsen (CG), \\
                      & Claas-Friso Hente (WiWi), Annkathrin Löffler, Andreas Hartkamp (Unicorns), \\
                      & Ajdina Karahasan (HWP), Ailina Salten, Lotte Rullkötter, \\
                      & Nabil Jabareen, Jan Detampel (MIN), Dominic Laumer (Bart), \\
                      & Karima Schulze, Henri Weber (AL) \\
                      &\\
        Unentschuldigt abwesend: & Tahnee Herzig (CG), Jakob Pape (BGZM),  Kay Zöllmer (WiWi), \\
                                & Samet Gunay (HWP), Johann Baumhoefener (Jura) \\
                                &\\
        Rücktritte: & Mirzo Ulugbek Khatamov \rightarrow Dean Haltmayer (CG)\\
    \end{tabular}
    \newpage
    %%% Vorgeschlagene Tagesordnung %%%
    \underline{Vorgeschlagene Tagesordnung}
    \begin{enumerate}[label={\textbf{Top \theenumi}},leftmargin=*]
        % folgende Zeile bei Sitzungsfortsetzungen einkommentieren und x durch
        % die Zahl des zuletzt behandelten TOPs ersetzen
        %\setcounter{enumi}{x}
        \item Vorbereitung der Wahlen zum Studierendenparlament (60 Min.)
        \begin{enumerate}
            \item Diskussion über Konzeptvorschlag des Präsidiums (V1617-034)
            \item Diskussion über die Steigerung der Wahlbeteiligung
            \item Beschlussfassung über den Zeitplan und den Urnenplan
            (V1617-035, V1617-036)
        \end{enumerate}
        \item Debatte über den Haushalt der Freien und Hansestadt Hamburg
              2017/2018 (V1617-038) (60 Min.)
        \item Aufruf zur Mahnwache zum 78. Jahrestag der Pogromnacht (V1617-037) (30 Min.)
        \item Erhöhung Wochenarbeitszeit für Festanstellung im Finanzreferat des AStA (V1617-039) (15 Min.)
        \item Kaffeeautomat im AStA-Infotreff (V1617-040) (15 Min.)
        \item Debatte zu G20 (V1617-041) (60 Min.)
        \item Verschiedenes (5 Min.)
    \end{enumerate}

    \newpage

    % Hier beginnt Teil A. Bei Fortsetzungen kommt hier nur der Teil über
    % die Feststellung der endgültigen Fassung der Tagesordnung.
    % Für sonstige Sitzungen ist dies bereits das Gerüst für die Formalia.
    % Man muss nur noch die protokollierten Inhalte eintragen.

    {\Large \underline{Teil A}}

    \subsection{Formalia}

    \textit{Jim eröffnet die Sitzung um 18.29 Uhr.}

    \subsubsection{Geschäftsbericht Präsidium}

    Jim berichtet von den Vorbereitungen der StuPa-Wahl für die Legislaturperiode
    2017/2018:

    \begin{itemize}
        \item Eine Sonderpostleitzahl ist eingerichtet.
        \item Hinsichtlich Druck und Gestaltung aller Dokumente sind erste Gespräche
              geführt und Angebote eingeholt worden.
        \item Die Formulare für die Kandidaturbögen sind veröffentlicht.
        \item Wie in den vergangenen Jahren sind fünf Stellen für Wahlhelfer*innen
              ausgeschrieben.
    \end{itemize}

    Ferner sei ein regelmäßiger E-Mail-Newsletter des StuPa-Präsidiums in
    Vorbereitung. Außerdem gebe es Bestrebungen die Internetseite des StuPa in
    Zukunft auch verschlüsselt abrufbar zu machen.

    In der kommenden Woche sei das Präsidium zusätzlich mit der Durchführung der
    Wahl zum FSR IGrLatPhil befasst.

    \subsubsection{Anfragen an das Präsidium}

    Keine

    \subsubsection{Geschäftsbericht AStA}

    Philipp berichtet von der "Welcome Week" des AStA und verweist darauf, dass
    viele interessante Gespräche währenddessen geführt worden seien. Hinsichtlich
    der Orientierungseinheiten der Fachbereiche bemängelt er, dass diese die
    "Feierei" teilweise als Selbstzweck sehen. Der AStA werde daher einen Dialog
    mit Fachbereichen und FSRen suchen, um sich über den Zweck der OEen
    auszutauschen.

    Franziska erklärt, dass die LAK ein gemeinsames Treffen mit den Leitern der
    Stadtteilschulen organisieren wolle. Hintergrund sei das im Sommer von den
    Schulleiterinnen und Schulleitern der Hamburger Stadtteilschulen beschlossene
    Positionspapier. Darüber hinaus sei der AStA aktuell in intensiven Gesprächen
    mit der Verwaltung der Universität bezüglich der Nutzungsmöglichkeiten der
    Foyers. Abschließend weist Franziska daraufhin, dass gegen Esther Bender
    keine öffentliche Klage wegen Untreue (Ausgaben der VS zugunsten einer
    Beteiligung am Christopher Street Day) erhoben werde und die Staatsanwaltschaft
    die Ermittlungen wegen Geringfügigkeit eingestellt habe.

    Philipp ergänzt seinen Bericht um den Hinweis, dass noch in diesem Jahr eine
    Veranstaltung mit Yanis Varoufakis durchgeführt werde.

    Timo erklärt, dass der Haushalt 2015/2016 abgeschlossen sei und der Haushalt
    2016/2017 erfolgreich vom Wirtschaftsrat genehmigt worden sei.

    Artur führt aus zum Thema "Wissenschaftsfreiheit statt Berufsverbote -
    ‚Verfassungsschutz‘ auflösen" (Anlage 1).  Aktuell erarbeite man überdies ein
    Thesenpapier zu den psychischen Leiden im Bologna-Studium.

    Armin erläutert, dass der Termin des G20-Gipfels in Hamburger näher rücke
    und man zurzeit die Einordnung des AStA innerhalb des Widerstandes diskutiere.
    Außerdem habe man sich mit dem Leistungsstipendium der Universität Hamburg
    befasst und konnte hierbei erreichen, dass künftig das soziale Engagement
    der Bewerberinnen und Bewerber stärker gewichtet werde.

    Oliver berichtet, dass seit dem letzten Zusammentritt des StuPa zwei
    "Stille Wörtchen" sowie der jährliche "KalendAStA" erschienen seien. Ferner
    sei das Öffentlichkeitsreferat bei der Vorbereitung von zwei Veranstaltungen
    involviert gewesen:

    \begin{itemize}
        \item "Gegen das Vergessen" - Veranstaltung des Auschwitz-Komitees zum
        Gedenken an die Pogromnacht 1938 (Anlage 2)
        \item Mahnwache anlässlich des 78. Jahrestages der Reichspogromnacht (Anlage 3)
    \end{itemize}

    Domenica berichtet über den Stand der Diskussion um die Rahmenprüfungsordnung.
    Derzeit werden in allen Teilen der Universitäten Stellungnahmen zum
    vorliegenden Entwurf erarbeitet. Ferner weist sie daraufhin, dass am
    30. November 2016 der nächste Dies Academicus stattfindet.

    Oliver ergänzt seinen Bericht dahingehend, dass auf Initiative des AStA
    Werbung von/für die Bundeswehr aus den Mensen des Studierendenwerks
    entfernt worden sei.

    Jakob weist darauf hin, dass ab sofort wieder alle zwei Wochen, und zwar im
    Wechsel zu den StuPa-Sitzungen, donnerstags eine Polit-Kneipe stattfinden
    werde; die jeweiligen Orte würden noch bekanntgegeben. Ferner starte nun die
    Veranstaltungsreihe "What´s left?". Hinsichtlich des Kulturkursprogrammes
    sei man zwar im Verzug, weil der Aufwand sehr groß sei, aber auch es werde
    demnächst starten.

    Golnar gibt einen Rückblick über die Themen und Veranstaltungen des Referats
    für internationale Studierende und weist daraufhin, dass demnächst eine
    Veranstaltung zum Friedensprozess in Kolumbien geplant sei. Überdies berichtet
    Golnar über die Auseinandersetzungen um das Studienkolleg:  Insbesondere
    kämpfe man darum, dass eine andere Form von Sprachtest angeboten werde.
    Künftig werde das Studienkolleg zudem nicht mehr zur Schulbehörde, sondern
    zur Wissenschaftsbehörde zählen. Dies habe jedoch zur Folge, dass es aktuell
    keine klaren rechtlichen Regelungen bezüglich des Studienkollegs gebe.

    \subsubsection{Anfragen an den AStA}

    Aufgrund vereinzelter Nachfragen wird eine allgemeine Diskussion über die
    Beteiligung der Verfassten Studierendenschaft an der Friedensbewegung geführt.

    Auf Nachfrage von Golnar erklärt Oliver, dass es aktuell keinen neuen Stand
    bezüglich eines Tages gegen Rechts gebe. Er spricht sich jedoch dafür aus,
    dass sich der Ausschuss gegen Rechts zeitnah treffen möge.

    \subsubsection{Dringlichkeitsanträge des AStA}

    Keine.

    \subsubsection{Aktuelle Stunde}

    Es liegt kein Antrag vor.

    \subsubsection{Feststellung der endgültigen Fassung des Teils B der Tagesordnung}

    Artur beantragt, dass die "Debatte zu G20" (V1617-041) [TOP-alt 6] vorgezogen
    werden soll und als TOP-neu 1 zu behandeln ist.

    Der Antrag wird mit zwölf Fürstimmen und zwölf Gegenstimmen, bei vier
    Enthaltungen abgelehnt. Aufgrund einer Anzweiflung wird die Abstimmung
    wiederholt: Dieses Mal entfielen 13 Fürstimmen und zwölf Gegenstimmen, bei
    zwei Enthaltungen, für den Antrag. Die Abstimmung wird erneut angezweifelt.

    Das Präsidium führt eine namentliche Abstimmung durch und verbindet dies mit
    der Feststellung der Beschlussfähigkeit (Anlage 4). Im Ergebnis wird der
    Antrag bei 14 Fürstimmen, elf Gegenstimmen und zwei Enthaltungen mehrheitlich
    angenommen.

    Die geänderte Tagesordnung wird bei einer Enthaltung einstimmig genehmigt.

    %%% BESCHLOSSENE TAGESORDNUNG %%%
    \underline{Endgültige Fassung der Tagesordnung}
    \begin{enumerate}[label={\textbf{Top \theenumi}},leftmargin=*]
        % folgende Zeile bei Sitzungsfortsetzungen einkommentieren und x durch
        % die Zahl des zuletzt behandelten TOPs ersetzen
        %\setcounter{enumi}{x}
        \item Debatte zu G20 (V1617-041) (60 Min.)
        \item Vorbereitung der Wahlen zum Studierendenparlament (60 Min.)
        \begin{enumerate}
            \item Diskussion über Konzeptvorschlag des Präsidiums (V1617-034)
            \item Diskussion über die Steigerung der Wahlbeteiligung
            \item Beschlussfassung über den Zeitplan und den Urnenplan
            (V1617-035, V1617-036)
        \end{enumerate}
        \item Debatte über den Haushalt der Freien und Hansestadt Hamburg
              2017/2018 (V1617-038) (60 Min.)
        \item Aufruf zur Mahnwache zum 78. Jahrestag der Pogromnacht (V1617-037) (30 Min.)
        \item Erhöhung Wochenarbeitszeit für Festanstellung im Finanzreferat des AStA (V1617-039) (15 Min.)
        \item Kaffeeautomat im AStA-Infotreff (V1617-040) (15 Min.)
        \item Verschiedenes (5 Min.)
    \end{enumerate}

    \subsubsection{Feststellung der Beschlussfähigkeit}

    \textit{Das StuPa ist mit 27 anwesenden Parlamentarier*innen beschlussfähig.}

    \subsubsection{Genehmigung der Protokolle der vorangegangenen Sitzungen}

    Das Protokoll der fünften Sitzung vom 16.06.2016 wird ohne Änderungen
    einstimmig genehmigt.

    Das Protokoll der sechsten Sitzung vom 30.06.2016 wird ohne Änderungen und
    bei drei Enthaltungen einstimmig genehmigt.

    Das Protokoll der siebten Sitzung vom 14.07.2016 wird ohne Änderungen
    einstimmig genehmigt.

    \textit{Präsidiumspause bis 20.24 Uhr}

    \vspace{0.5cm}
    {\Large \underline{Teil B}}

    \subsection{Debatte zu G20}

    Armin stellt den Antrag für die Antragsteller*innen vor.

    Benjamin zweifelt die Beschlussfähigkeit des Studierendenparlamentes an.

    \textit{Die Beschlussfähigkeit ist mit 19 von 47 Parlamentarier*innen nicht
    gegeben.}

    \textit{Jim erklärt die Sitzung um 20.39 Uhr für abgebrochen.}
\end{document}
